\addtoindex{Notte maledetta sei per me...}


\begin{verse}
    Notte maledetta sei per me\\
    la relazione più tossica.\\!
    Ti voglio ti bramo\\
    e quando sei qui per me\\
    non mi concedo.\\!

Mi devo svuotare %
\tikz[remember picture,baseline=(punto.base)]%
  \node[inner sep=0pt,outer sep=0pt] (punto) {per};%
\end{verse}

\begin{tikzpicture}[remember picture,overlay]
  \pgfmathsetseed{1234}

  % punto di partenza leggermente a destra
  \coordinate (partenza) at ([xshift=-0.3pt, yshift=-0.8pt]punto.north east);

  \begin{scope}[shift={(partenza)}]
    % curva verso destra e in basso
    \draw[thick, black, smooth, decorate,
          decoration={random steps, segment length=4pt, amplitude=0.5pt}]
      (0,0) .. controls (0.4,-0.1) and (1,-1) .. (0.7,-4); 
  \end{scope}
\end{tikzpicture}


  % \begin{scope}[shift={(partenza)}]
  %   \draw[thick, black!70, smooth, decorate,
  %         decoration={random steps, segment length=4pt, amplitude=01pt}]
  %     (0,0) -- (0.5,-2cm);




  % Ora disegno la linea tremolante
  % \begin{scope}[shift={(partenza)}]
  %   \draw[thick, black!70]
  %     plot[
  %       domain=0:-8cm,
  %       samples=50,
  %       smooth
  %     ]
  %     ({0.1*(rand)}, {\x});
  % \end{scope}
% \tikz[remember picture,baseline=(punto.base)]%
%   \node[inner sep=0pt,outer sep=0pt] (punto) {per};%
% \end{verse}

% % linea tremolante overlay
% \begin{tikzpicture}[remember picture,overlay]
%   \pgfmathsetseed{1234} % opzionale: rende il "caso" riproducibile

%   % SHIFT: metto l'origine (0,0) esattamente in corrispondenza del punto di partenza
%   \begin{scope}[shift={((punto.north east)+(2pt,0)}]
%     % ora disegno la linea come plot verticale: x = piccolo rumore, y = valore di \x (dominio)
%     \draw[thick, black!70]
%       plot[
%         domain=0:-8cm,   % scende di 8cm
%         samples=50,
%         smooth
%       ]
%       ({0.1*(rand)}, {\x}); % x = tremolio relativo, y = \x (dominio)
%   \end{scope}
% \end{tikzpicture}

% \begin{verse}
%     Notte maledetta sei per me\\
%     la relazione più tossica.\\!
%     Ti voglio ti bramo\\
%     e quando sei qui per me\\
%     non mi concedo.\\!

% Mi devo svuotare per \tikz[remember picture,baseline=(punto.base)]%
%   \node[inner sep=0pt,outer sep=0pt] (punto) {per};%
% \end{verse}

% % il tikzpicture overlay va dopo il testo (può essere messo anche a fine pagina)
% \begin{tikzpicture}[remember picture,overlay]
%   \draw[thick, black!70] plot[domain=0: -8, samples=40, smooth] 
%     ({2 + 0.1 * (rand)}, {\x});
%   % \draw[
%   %   decorate,
%   %   decoration={random steps,segment length=1pt, amplitude=1pt},
%   %   thick, color=black!70
%   % ] (punto.south) .. controls +(0,-3mm) and +(0,3mm) .. ++(0,-8cm);
% \end{tikzpicture}

% %     Mi devo svuotare per
    
% % \tikz[remember picture, baseline] \node (puntoPartenza) {svuotare};

% % \begin{tikzpicture}[remember picture, overlay]
% %   \draw[
% %     decorate,
% %     decoration={random steps, segment length=1pt, amplitude=1pt},
% %     thick, color=black!70
% %   ] (puntoPartenza.south) -- ++(0,-8cm);
% % \end{tikzpicture}
% % \begin{tikzpicture}
% %   \pgfmathsetseed{1234} % per rendere il “caso” riproducibile

% %   \draw[thick, black!70] plot[domain=0: -8, samples=40, smooth] 
% %     ({2 + 0.1 * (rand)}, {\x}); % aggiunge rumore su x

% %   % se vuoi rumore solo su y, puoi fare qualcosa come
% %   % plot[domain=0:-8, samples=40, smooth] ({2}, {\x + 0.1 * rand})
% % \end{tikzpicture}

    
%     % \begin{tikzpicture}[decoration={bent, amplitude=3pt, aspect=0.4}]
%     %   \draw[decorate, thick] (0,0) -- (3,1) -- (1,2);
%     % \end{tikzpicture}

%     % \tikz[baseline] \node[draw=none, pencildraw/.style={
%     %     black!75,
%     %     decorate,
%     %     decoration={random steps, segment length=0.8pt, amplitude=0.1pt}
%     %   }
    


%     % \begin{tikzpicture}[pencildraw/.style={
%     %     black!75,
%     %     decorate,
%     %     decoration={random steps, segment length=0.8pt, amplitude=0.1pt}
%     %   }]
%     %   \draw[pencildraw, thick] (0,0) -- (5,0) to[out=30, in=330] (7,1);
%     % \end{tikzpicture}




%     % \tikz[baseline]{
%     %     \draw[decorate, decoration={random steps, segment length=3pt, amplitude=3pt}]
%     %         (0,0) -- ++(0,-18cm);
%     % }

% \end{verse}